\documentclass[11pt]{article}
\usepackage[a4paper,margin=1in]{geometry}
\usepackage[T1]{fontenc}
\usepackage[utf8]{inputenc}
\usepackage{setspace}
\usepackage{hyperref}

\setstretch{1.1}
\setlength{\parskip}{0.6em}
\setlength{\parindent}{0pt}

\begin{document}

\begin{flushright}
May 16, 2026
\end{flushright}

Editors \\
\textit{IEEE Access}

Dear Editors,

Please consider our manuscript entitled

\textbf{``When Agreement Fails: Visual Anchoring and Multimodal Pseudo-Label Reliability in Classroom Behavior Recognition''}

for publication in \textit{IEEE Access}.

We submit this paper as a reliability-audit study of multimodal large
language model (MLLM) pseudo-labeling in classroom behavior recognition.
Using three SCB sub-datasets, we organise the evidence into four
linked analysis modules: a protocol-matched zero-shot CLIP baseline, MLLM
annotation and cross-model validation, CLIP adaptation under label-source
choices, and mechanism/strategy audits. The study addresses a practical
deployment question that is increasingly important across computer vision and
educational AI: when can general-purpose MLLMs provide useful pseudo-label
supervision, and when does cross-model agreement simply preserve shared error?

The manuscript makes four main contributions. First, it provides a
reproducible audit protocol that links annotation quality, downstream
adaptation, retention controls, distribution-shift checks, anchor-proxy
analysis, auxiliary filtering and self-training baselines, and cross-model
robustness. Second, it shows that pseudo-label usefulness depends on
category-level visual evidence. Third, it demonstrates that agreement
filtering is not a portable denoising rule: in the observed setting,
agreement-filtered subsets can be affected by data-volume loss,
class-composition shift, and shared label error. Fourth, it separates these
boundary conditions from a single annotator choice or optimisation path using
repeated-seed downstream evaluation, targeted LoRA replications, train-split
teacher--student and confidence-filtering baselines, and a five-model
robustness check.

We believe the manuscript is well suited to \textit{IEEE Access} because it
offers a reproducible, application-driven study at the intersection of
multimodal AI, computer vision, and intelligent education. The emphasis is not
another benchmark comparison; it is an auditable methodology and evidence base
for designing reliable MLLM-based annotation workflows in classroom sensing and
related vision systems.

This submission is original, has not been published previously, and is not under
consideration elsewhere. All authors have approved the submitted version and
agree to its submission to \textit{IEEE Access}. The authors declare no
competing interests.

Funding statement: This research was funded by the Education Department of
Hunan Province, grant number HNJG-2022-0608.

Data and code availability: The SCB dataset is publicly available at
\url{https://huggingface.co/datasets/wintonYF/SCB-Dataset}. The experiment code
is available at \url{https://github.com/zhanglizhuo/llm-annotation}.

Ethics and privacy: This study is a secondary analysis of a publicly released
benchmark. The manuscript reports aggregate results, makes no re-identification
attempts, and uses privacy-preserving pixelated representative crops where
visual examples are needed.

Thank you for your time and consideration.

Sincerely,

Yan Ma and Lizhuo Zhang

Corresponding author:
Lizhuo Zhang \\
School of Information and Intelligent Science and Technology, \\
Hunan Agricultural University, Changsha 410128, China \\
Email: zhanglizhuo@hunau.edu.cn

\end{document}